L'experiment de Millikan va permetre determinar la càrrega de l'electró. El muntatge d'aquest experiment consta de dues plaques metàl·liques horitzontals, una sobre de l'altra, separades verticalment per una distància $d$ i connectades a una font de potencial elèctric regulable. En l'espai que hi ha entre les plaques s'introdueixen algunes gotetes d'oli carregades negativament. L'experiment consisteix a crear un camp elèctric entre les plaques i aconseguir una posició d'equilibri de les gotetes d'oli contrarestant el seu pes.

\begin{apartat}{1,25}
Feu un esquema del dispositiu emprat per Millikan, dibuixant les forces que actuen sobre una gota d'oli esfèrica. Indiqueu i raoneu el signe de la càrrega que tindrà cada placa i la direcció i el sentit del camp elèctric generat per aquestes plaques.\\
\hspace*{1em}Suposant que les plaques es connecten a un potencial elèctric de 2,00~kV i que la distància entre plaques és $d = 2,00\ \mathrm{cm}$, calculeu el camp elèctric creat entre plaques i dibuixeu-lo en el mateix esquema.
\end{apartat}

\begin{apartat}{1,25}
Tenint en compte que la densitat de l'oli és de $923\ \mathrm{kg\,m^{-3}}$ i que el radi d'una gota és d'1,08~$\mu$m, calculeu la càrrega d'una gota que es troba en equilibri. Quants electrons calen per a generar aquesta càrrega?\\
\hspace*{1em}Què s'observaria si il·luminéssim la gota amb raigs ultraviolats i aquesta perdés un electró? Justifiqueu la resposta.
\end{apartat}

\dades{%
  $|e| = 1,602\times 10^{-19}\ \mathrm{C}$.\\
  $m_\mathrm{e} = 9,11\times 10^{-31}\ \mathrm{kg}$.\\
  $g = 9,8\ \mathrm{m\,s^{-2}}$.\\
  Volum d'una esfera: $\dfrac{4}{3}\pi r^3$.}
